\section*{1. Mục tiêu và kiến trúc hệ thống}
\addcontentsline{toc}{section}{1. Mục tiêu và kiến trúc hệ thống}

\subsection*{1.1. Yêu cầu bài thực hành}

Bài thực hành xây dựng một giao diện tương tác người--máy trong đó biểu cảm
khuôn mặt quan sát qua webcam được chuyển thành lệnh di chuyển cho
TurtleBot~4. Sáu trạng thái cần xử lý là \textit{happy}, \textit{sad},
\textit{neutral}, \textit{angry}, \textit{surprise} và \textit{fear}. Mỗi
trạng thái được ánh xạ tới một hành vi đơn giản: cong trái, cong phải, dừng,
lùi, xoay tròn hoặc tiến.

Theo yêu cầu vận hành cuối cùng, hệ thống xử lý ngay observation đầu tiên khi
có đúng một khuôn mặt và confidence lớn hơn 0,60; không dùng cửa sổ ổn định
theo thời gian. Mỗi kích hoạt vẫn chỉ được thực hiện một lần; mất khuôn mặt,
mất kết quả inference hoặc mất odometry phải đưa vận tốc về không. Chế độ
\texttt{preview\_only=true} được bật mặc định để kiểm tra camera và logic mà
không làm robot chuyển động.

Package được phát triển độc lập với tên \texttt{tbe4\_face\_recognition}, đặt
tại \path{~/Interaction/tbe4_face_recognition}. Nền tảng mục tiêu là
Ubuntu~22.04, Python~3.10 và ROS~2 Humble \cite{ros2humble}; phần mô phỏng sử
dụng TurtleBot~4 Standard và simulator chính thức \cite{turtlebot4}.

\subsection*{1.2. Cấu trúc package}

\begin{svcode}
tbe4_face_recognition/
|-- config/
|   |-- recognition.yaml
|   `-- controller.yaml
|-- launch/
|   |-- recognition.launch.py
|   |-- interaction.launch.py
|   `-- simulation.launch.py
|-- scripts/
|   |-- setup_face_venv.sh
|   |-- build.sh
|   |-- run_recognition.sh
|   `-- run_interaction.sh
|-- tbe4_face_recognition/
|   |-- protocol.py, emotion_model.py
|   |-- emotion_filter.py, behavior.py, motion.py
|   |-- emotion_recognition_node.py
|   `-- emotion_controller_node.py
|-- test/
|-- package.xml, setup.py, setup.cfg
`-- README.md
\end{svcode}

Các module \texttt{emotion\_filter.py}, \texttt{behavior.py} và
\texttt{motion.py} là logic Python thuần, không phụ thuộc TensorFlow hay ROS.
Hai node chỉ đảm nhiệm I/O, lifecycle và logging. Cách chia này cho phép kiểm
thử phần quyết định và chuyển động bằng mô phỏng số mà không cần webcam hoặc
Gazebo.

\subsection*{1.3. Kiến trúc hai node}

\begin{figure}[H]
\centering
\resizebox{\textwidth}{!}{%
\begin{tikzpicture}[
  box/.style={draw,rounded corners,align=center,minimum height=1.15cm,
              text width=2.8cm,fill=blue!5,font=\scriptsize},
  topic/.style={draw,align=center,minimum height=1.0cm,
                text width=2.55cm,fill=orange!10,font=\scriptsize},
  arrow/.style={-{Latex[length=2.5mm]},thick}
]
  \node[box] (camera) at (0,0) {Webcam laptop\\OpenCV VideoCapture};
  \node[box] (recognition) at (3.9,0)
    {emotion\_recognition\_\\node\\FER/TFLite, overlay};
  \node[topic] (observation) at (7.8,0)
    {/emotion\_\\observation\\std\_msgs/String JSON};
  \node[box] (controller) at (11.7,0)
    {emotion\_controller\_\\node\\filter, state machine};
  \node[topic] (cmd) at (15.6,0)
    {/cmd\_vel\\geometry\_msgs/\\Twist};
  \node[box] (robot) at (19.5,0) {TurtleBot 4\\motion control};
  \node[topic] (status) at (7.8,-2.5)
    {/emotion\_behavior\_\\status\\std\_msgs/String JSON};
  \node[topic] (odom) at (15.6,-2.5)
    {/odom\\nav\_msgs/\\Odometry};

  \draw[arrow] (camera) -- (recognition);
  \draw[arrow] (recognition) -- (observation);
  \draw[arrow] (observation) -- (controller);
  \draw[arrow] (controller) -- (cmd);
  \draw[arrow] (cmd) -- (robot);
  \draw[arrow] (robot.south) |- (odom.east);
  \draw[arrow] (odom.west) -| (controller.south);
  \draw[arrow] (controller.south west) |- (status.east);
  \draw[arrow] (status.west) -| (recognition.south);
\end{tikzpicture}
}
\caption{Kiến trúc dữ liệu của hệ thống nhận diện biểu cảm và điều khiển robot.}
\end{figure}

Node nhận diện chạy bằng Python trong virtualenv riêng, sở hữu webcam và cửa
sổ OpenCV. Node điều khiển chạy bằng Python hệ thống của ROS, không import FER
hoặc TensorFlow, và là publisher duy nhất của \path{/cmd_vel}. Vì hai node là
hai process độc lập, một lần inference chậm không chặn timer điều khiển 20~Hz.
Nếu observation không được cập nhật, watchdog của controller sẽ dừng robot.

\subsection*{1.4. Topic và contract dữ liệu}

\begin{longtable}{@{}p{0.29\textwidth}p{0.67\textwidth}@{}}
\toprule
\textbf{Topic} & \textbf{Kiểu dữ liệu và vai trò} \\ \midrule
\endfirsthead
\toprule \textbf{Topic} & \textbf{Kiểu dữ liệu và vai trò} \\ \midrule
\endhead
\path{/emotion_observation} &
\path{std_msgs/msg/String}; mang JSON gồm session, sequence, nhãn FER,
confidence và số khuôn mặt. \\
\path{/emotion_behavior_status} &
\path{std_msgs/msg/String}; phản hồi state, biểu cảm ổn định, hành động đang
chạy, reason, trạng thái preview/odometry và tiến độ. \\
\path{/odom} &
\path{nav_msgs/msg/Odometry}; pose tương đối dùng để đo quãng đường và góc quay. \\
\path{/cmd_vel} &
\path{geometry_msgs/msg/Twist}; vận tốc thẳng và vận tốc góc gửi tới robot. \\
\bottomrule
\end{longtable}

Một observation hợp lệ có dạng:

\begin{svcode}
{"confidence":0.91,"emotion":"happy","face_count":1,
 "sequence":42,"session_id":"a1b2c3d4"}
\end{svcode}

\texttt{session\_id} thay đổi khi node nhận diện khởi động lại; sequence tăng
sau mỗi inference hoàn thành. Controller chỉ coi một sequence mới trong cùng
session là heartbeat hợp lệ. Protocol kiểm tra đúng tập khóa, kiểu dữ liệu,
miền confidence và nhãn; payload sai không được truyền thẳng vào state machine.
